\documentclass[conference]{IEEEtran}
\usepackage{cite}
\usepackage{amsmath,amssymb}
\usepackage{graphicx}
\begin{document}
\title{From Soil Mechanics to Crop Stress: Screening Regolith Simulants for Plant Suitability from Public Soil Data}
\author{\IEEEauthorblockN{Ikon Kim}
\IEEEauthorblockA{High School Student\\St. Mark's School}}
\maketitle
\begin{abstract}
Growing crops on the Moon or Mars means growing them in regolith, and testing whether plants tolerate a given simulant takes a multi-week growth trial, so only a few of the many published simulants have been tested. This work asks whether public soil measurements alone can screen simulants for how hard they are on plants, before any trial. Three public datasets, engineering properties, grain size, and plant results, are joined into one table keyed by simulant. A random forest that tries to predict cohesion from the other properties fails under mission-grouped cross validation (R2 below zero); this negative result is reported because it reflects how noisy the pooled cohesion data is. In its place a screening score is built from measured cohesion and a small pH-to-stress calibration, and 23 simulants are ranked from friendliest to riskiest. The ranking agrees with the two simulants that have published growth results, and a jitter test shows the friendly and risky ends are stable while the middle is not. The result is a reproducible screen that flags likely problem soils from data that already exists.
\end{abstract}
\begin{IEEEkeywords}
data integration, screening index, regolith simulants, space agriculture, ranking, reproducibility
\end{IEEEkeywords}
\section{Introduction}
Any long stay on the Moon or Mars will eventually need to grow food in the local regolith, because shipping soil from Earth does not scale. Regolith is a poor medium: the material tested so far is alkaline, holds few nutrients, and packs into a hard surface crust that keeps water from reaching roots \cite{russell}. To prepare, labs make simulants, ground rock mixtures meant to stand in for real regolith, and there are now dozens of them. Testing whether plants grow in a simulant means a growth trial that runs for weeks, so only a few have ever been grown in. There is no quick way to guess, up front, which of the many simulants are hardest on plants.
The numbers needed for a first guess may already exist, just in the wrong place. Soil engineers have measured cohesion, bulk density, and grain size for many of these simulants \cite{gasteiner,planetgsd}, and those properties decide whether a soil compacts and crusts. Plant biologists, separately, have measured pH and growth for a few of them \cite{russell}. This project brings the two together into one screening score, so a simulant can be flagged as likely friendly or likely harsh from published measurements alone, before anyone grows a single plant in it.
The paper makes four contributions. It joins three public datasets into one table keyed by simulant. It tests whether a model can predict cohesion from the other properties, and reports honestly that it cannot under grouped cross validation. It builds a screening score from measured cohesion plus a small pH calibration and ranks 23 simulants. And it checks that ranking two ways, against the simulants with published growth results and with a jitter test for stability. The negative model result is kept as part of the story, not cut to make the work look cleaner.
\section{Related Work}
Several groups have put plants in regolith simulant. Wamelink and coworkers \cite{wamelink} germinated seeds in diluted Mars and Moon simulants, though the plants stayed small. Eichler \cite{eichler} found that on plain Martian simulant with no amendment almost nothing grew. Most relevant here, Russell \cite{russell} grew lettuce, radish, and pepper in a carbonaceous asteroid simulant mixed with peat and watched growth fall as the simulant fraction rose, with radish worst hit and compaction and crusting blamed over missing nutrients. A recent review \cite{duri} adds that lunar simulants in particular are alkaline and nutrient poor, with JSC-1A giving poor growth and LHS-1 managing a little.
On the engineering side there is measured data that is rarely used for plant questions. Gasteiner and coworkers \cite{gasteiner} compiled an open database of lunar regolith and simulant properties, with cohesion, friction angle, and bulk density across many missions. PlanetGSD \cite{planetgsd} does the same for grain size, and others have tried to relate cohesion to bulk density for compacted simulants \cite{dotson}. These numbers describe whether a soil compacts and crusts, but they sit in a separate literature from the plant work.
So there are two piles of data about the same soils that almost never appear together. Screening indices are common in soil and materials science, but none was found that combines engineering and biology measurements to pick a soil for space crops. Treating this as a data integration problem \cite{doan}, and being explicit about which parts of the model work and which do not \cite{gundersen}, is the angle taken here.
\section{Data}
Three public datasets are used. The Gasteiner open database \cite{gasteiner} has measured cohesion, friction angle, and bulk density for lunar regolith and simulants across many missions and lab studies. PlanetGSD \cite{planetgsd} gives grain-size curves, and the OSD-670 study \cite{russell} supplies the plant side, with soil pH and cation exchange capacity for each peat and simulant mixture. Nothing here is an original experiment; all of it is reused public data, with sources and licenses logged in the repository.
The cohesion values needed cleaning first, because older reports give ranges such as "0.15 - 0.7" rather than a single number, so the midpoint of each range is taken. From the sieve curves, d10, d50, and d90 are computed per sample by sorting sizes, normalizing the weights, and interpolating the cumulative passing curve. After joining, 23 simulants carry a usable measured cohesion and enter the ranking.
\section{Method}
The first idea was to predict cohesion from grain size and bulk density with a random forest, to fill in simulants with no measured cohesion. To keep the test honest, cross validation is grouped by mission, so the model never trains on a soil from the same body it is tested on. It does not work: the grouped-CV R2 comes out below zero, worse than predicting the mean. Cohesion barely correlates with density here (Spearman about 0.13). The likely cause is that the cohesion numbers come from very different tests, from spacecraft landing estimates to lab shear boxes, so they do not fall on one clean curve. This is reported rather than hidden.
Since the model does not work, the measured cohesion is used directly for the simulants that have it. Higher cohesion means a soil holds together harder and is more likely to crust, so a rescaled cohesion is treated as a crusting risk in [0, 1]. One simulant, NAO-1, sits far above the rest near 95 kPa, so the scale is capped at the 95th percentile before rescaling, otherwise it flattens everything else into one narrow band.
For the chemistry side the OSD-670 data \cite{russell} is used. As the simulant fraction rises, soil pH rises and radish biomass drops. The drop is measured as a deficit against the all-peat control and fit against pH with a straight line, giving deficit = 0.351 pH - 1.675. The fit is tight (R2 = 0.975) but rests on only four points, so it is treated as a rough calibration, not a validated model.
The index combines the two. Where a simulant has a measured pH, the score is 0.4 times the crusting risk plus 0.6 times the calibrated chemistry stress; where it does not, and most do not, the score falls back to the crusting risk alone. In practice only one ranked simulant (JSC-1A) has a published pH, so the chemistry term moves that single row and the ranking is otherwise driven by compaction. The two sides do not yet carry equal weight, and the paper does not pretend they do.
\section{Results}
The index runs over the 23 simulants with a measured cohesion. Each gets a score in [0, 1], where 0 is friendliest to plants and 1 is riskiest. The scores are not evenly spread: most simulants sit below about 0.5 and only a handful reach the risky end, so the ranking is about picking out the few problem soils rather than splitting a smooth gradient. At the friendly end are LSS-ISAC-1 (0.00), EAC-1A (0.00), and MLS-1 (0.02), with LHS-1 (0.07) close behind. At the risky end IGG-01 and NAO-1 reach 1.00, with OB-1A (0.79) and PolyU-1 (0.71) just below; NAO-1 carries the very high measured cohesion and was always going to land here.
JSC-1A is the informative case, as the only simulant in the table with a published pH and therefore the only score carrying the chemistry term. Its compaction alone is low (0.08), but the high pH raises the chemistry stress and the combined score reaches 0.63, moving it into the risky half. This matches the plant studies, and shows how much one row can move once its pH is known.
Two simulants have published plant results to check against. JSC-1A, which the index places in the risky half at rank 19, is reported to cut germination and biomass \cite{duri}. LHS-1, placed in the friendlier half at rank 7, does sustain limited growth in the same line of work. On the two soils with plant data the direction of the ranking agrees. Two points is a sanity check, not a validation, but it is reassuring that the ranking does not point the wrong way.
To test stability, every cohesion value was jittered by up to 20 percent and the whole index rebuilt 2000 times. The friendliest three simulants stayed in the friendly group in about 78 percent of runs and the riskiest three in about 67 percent, with a median Kendall tau of 0.95 against the original order. The order barely moves overall; where it moves is the middle, where several simulants sit within a few hundredths and swap places. So the two ends are solid and the middle band should not be read as an exact order.
It is worth being clear about what the ranking is. For 22 of the 23 simulants the score is driven entirely by compaction, because they have no published pH. So this is mostly a physical, soil-structure ranking with a chemistry correction that so far touches one row. It is a screening tool, a way to point at the simulants most likely to crust and stress plants before anyone spends months on a growth trial, not a prediction of exact biomass.
\section{Discussion}
The main point is that a useful first pass at which simulants will be hard to grow in is possible without running a single growth trial. The measured properties already exist and carry enough signal to separate loose, friendly soils from dense, crust-prone ones. Because trials are slow and expensive and there are far more candidate simulants than any lab can test, a screen that puts likely problem soils at the top lets people spend trial time where it counts. Notably, the simple physical measure did most of the work; the chemistry probably matters more in reality, but compaction is what public data let me measure across the whole set, which is itself a finding.
The limits are worth naming. The chemistry calibration rests on four points from one study, so the pH-to-stress line is a rough guide. Only one ranked simulant has a published pH, so the chemistry term hardly touches the ranking. The cohesion values come from very different tests, which is exactly why the random forest could not fit them. And the plant check covers only two soils. None of these break the screen, but they set a ceiling on how much weight the exact scores can carry. The clearest next step is more pH and CEC values, so the chemistry side does real work across the table rather than on a single row.
\section{Conclusion}
This work asked whether public soil data alone can flag which regolith simulants are likely to crust and stress plants, without growing anything. It can, at least for the two ends of the ranking. A rescaled cohesion picks out the soils that pack down hardest, a four-point pH calibration adds a chemistry correction where the data allows, and a jitter test shows the friendly and risky groups are stable while the middle is not. The negative modeling result is reported alongside the positive screen. For a first pass at a hard problem, pointing at the likely trouble soils before anyone plants a seed is already useful.
\begin{thebibliography}{9}
\bibitem{russell} R. Russell et al., "Testing CI asteroid regolith simulant as a plant growth medium," Planetary Science Journal, vol. 3, no. 7, art. 155, 2022.
\bibitem{gasteiner} L. Gasteiner, N. Murdoch, and A. D'Angelo, "An open database of lunar regolith and simulant mechanical properties," arXiv:2602.03829, 2026.
\bibitem{planetgsd} PlanetGSD 1.0: a global grain-size database for planetary surfaces, Earth System Science Data, 2026.
\bibitem{wamelink} G. W. W. Wamelink et al., "Can plants grow on Mars and the Moon: a growth experiment on Mars and Moon soil simulants," PLOS ONE, vol. 9, no. 8, e103138, 2014.
\bibitem{eichler} A. Eichler et al., "Challenging the agricultural viability of Martian regolith simulants," Icarus, vol. 354, art. 114022, 2021.
\bibitem{duri} L. G. Duri et al., "The potential of regolith simulants as agricultural substrates," Frontiers in Astronomy and Space Sciences, vol. 8, art. 747821, 2022.
\bibitem{dotson} J. Dotson et al., "Cohesion and bulk density relationships for compacted regolith," Icarus, vol. 411, art. 115943, 2024.
\bibitem{doan} A. Doan, A. Halevy, and Z. Ives, Principles of Data Integration. Morgan Kaufmann, 2012.
\bibitem{gundersen} O. E. Gundersen and S. Kjensmo, "State of the art: reproducibility in artificial intelligence," in Proc. AAAI, 2018.
\end{thebibliography}
\end{document}
